\begin{abstract}
Large language models (LLMs) inherently absorb harmful knowledge, misinformation, and personal data during pretraining on large-scale web corpora, with no native mechanism for selective removal. While machine unlearning offers a principled solution, existing approaches are provider-centric, requiring retraining pipelines, curated retain datasets, and direct intervention by model service providers (MSPs), thereby excluding end users from controlling their own data. We introduce \textit{Interactive Machine Unlearning} (IMU), an inference-time setting in which users instruct LLMs to forget targeted knowledge through natural language, and which imposes two constraints that no existing unlearning method satisfies simultaneously, namely \textit{training-free} and \textit{single-sample} operation. To realize IMU, we propose \textbf{RePAIR}, a prompt-aware model repair framework comprising (i) a watchdog model for unlearning intent detection, (ii) a surgeon model for generating repair procedures, and (iii) a patient model whose parameters are updated autonomously. At the core of RePAIR, we develop \textbf{S}teering \textbf{T}hrough \textbf{A}ctivation \textbf{M}anipulation with \textbf{P}seudo\textbf{I}nverse \textbf{(STAMP)}, a training-free, single-sample unlearning method that redirects MLP activations toward a refusal subspace via closed-form pseudoinverse updates. Its low-rank variant reduces computational complexity from $O(m^3)$ to $O(r^3 + r^2 \cdot m)$, where $m$ is the feed-forward width, enabling efficient on-device unlearning with up to $\sim$3$\times$ speedup over training-based baselines. Extensive experiments across harmful knowledge suppression, misinformation correction, and personal data erasure demonstrate that RePAIR achieves near-zero forget scores ($\mathrm{Acc}_f = 0.00$, $F\text{-}RL = 0.00$) while preserving model utility ($\mathrm{Acc}_r$ up to 84.47, $R\text{-}RL$ up to 0.88), outperforming six state-of-the-art baselines. These results establish RePAIR as an effective and practical framework for user-driven model editing, advancing transparent and on-device control over learned knowledge, with potential extensions to multimodal foundation models.
\end{abstract}